% =========================================================================
% Metaheuristics Project (31-MM34, SoSe 26) -- Report Structure (V1)
% Skeleton only: sections, subsections and one-line content notes.
% The full text lives in Report_Full_V1.tex (identical structure).
% Compile: pdflatex Report_Structure_V1.tex
% =========================================================================
\documentclass[10pt, a4paper]{article}

\usepackage[utf8]{inputenc}
\usepackage[T1]{fontenc}
\usepackage{lmodern}
\usepackage[margin=2.5cm]{geometry}
\usepackage{amsmath,amssymb}
\usepackage{booktabs}
\usepackage[hidelinks]{hyperref}

\title{An Adaptive Large Neighborhood Search for the\\
       Prize-Collecting Skill Vehicle Routing Problem}
\author{Umais Chaudhary, Rehman Kerimov, Dustin Schulz, Marco Leander vom Orde \\
        \small 31-MM34 Metaheuristics, SoSe 26, Bielefeld University}
\date{July 2026}

\begin{document}
\maketitle

\begin{abstract}
% ~130 words. Problem (selective routing + skills + time windows + shifts,
% short CPython budget); development path: literature -> ALNS; tuned LNS
% baseline as yardstick (63.0% UB at 30 s); modular ALNS with 14 destroy /
% 9 repair operators, penalized SA acceptance, time-based cooling,
% runtime-normalized adaptive weights; grid search + leave-one-out ablation
% (every operator contributes); headline results (matches/beats baseline at
% equal budget; 65.5% UB average in 10-minute runs, all runs feasible).
\end{abstract}

% Mapping to the four components required by the project description
% (31-MM34, "Submission document"):
%   1. Approach description                      -> Sections 2 and 4
%   2. Citations / discussion of related work    -> Section 3
%   3. Experimental investigation: components    -> Section 5  (!!!)
%   4. Experimental investigation: overall perf. -> Section 6  (!!!)
% No code in the document; pseudocode (Algorithm 1) in course-slides style.

% -------------------------------------------------------------------------
\section{Introduction}
% Motivation (pharmacy home delivery). Problem features: optional customers,
% skills, time windows, shifts; hard wall-clock budget in pure Python.
% Red thread of the report = development path: (i) evidence from the
% literature fixes the method family, (ii) a strong-but-simple LNS baseline
% fixes the yardstick, (iii) engineering makes evaluation cheap, (iv) breadth
% of operators is added and then *measured* (grid search, ablation, weight
% tracking) rather than argued. Contributions list.

% -------------------------------------------------------------------------
\section{Problem Statement}
% Formal definition: customers C, couriers B; profits, service times, time
% windows, shifts, skill sets; rounded Euclidean travel times; objective;
% feasibility rules. Singleton upper bound UB as evaluation yardstick and
% its limitations.

% -------------------------------------------------------------------------
\section{Related Work and Choice of Method}
\subsection{Problem classification}
% Orienteering family / TOPTW at the routing core; Skill VRP and technician
% routing (TRSP) for the skill aspect; NP-hardness.
\subsection{Metaheuristics for selective routing with time windows}
% ILS (Vansteenwegen), granular VNS (Labadie), I3CH (Hu & Lim), tuned ILS/SA
% (Gunawan); forward time slack as the shared efficiency idea.
\subsection{LNS, ALNS and population methods}
% Shaw; Ropke & Pisinger; SISR; ALNS for the TRSP (Kovacs); HGS (Vidal) and
% why it fits poorly here (customer selection + hard feasibility filters).
\subsection{Synthesis}
% Both closest problem families point to destroy-and-repair search with
% insertion-based repair and annealing acceptance -> ALNS; reinforced by the
% short interpreted-Python budget.

% -------------------------------------------------------------------------
\section{Approach Description}
\subsection{Development path and baseline LNS}
% Milestone-driven development. Baseline solver skillvrp+.py: multi-start
% greedy insertion + LNS with record-to-record acceptance; iteratively tuned
% from 61.8 to 63.0% UB (30 s benchmark logs). Role: independent yardstick
% so every ALNS design decision is measured, not argued.
\subsection{Framework overview}
% Algorithm 1 (package-free tabbing figure): weighted operator selection,
% penalized SA acceptance, time-based cooling, segment-wise weight updates,
% restart + reheat after stagnation, optional polish hook.
\subsection{Preprocessing and construction heuristic}
% Distance matrix; per-customer solo-feasible courier lists (skill filter +
% singleton tour check); profit-density bound. Multi-start greedy insertion
% with four orderings (hybrid / profit / density / fewest couriers).
\subsection{Destroy operators}
% 14 operators in five groups: random; cost-based (worst-density,
% worst-detour v1/v2, largest-saving); relatedness-based (related, Shaw,
% temporal Shaw, cluster); structural (route, sequence, time-window
% segment); history-based; problem-specific skill-scarcity removal.
% Destroy size sampled per call from fraction ranges.
\subsection{Repair operators}
% 9 operators: greedy best-first, regret-2, regret-3 (courier-level),
% sequential cheapest (profit order), noised greedy, scarce-skill-first,
% last-removed-first-inserted, random-position, Shaw insertion.
% Candidate set = removed customers + capped sample of unserved extras.
\subsection{Efficient move evaluation}
% Route caches with forward time slacks -> O(1) insertion feasibility;
% per-(customer, vehicle) move cache with per-route invalidation; penalized
% evaluation (time-window/shift penalties, prohibitive skill penalty).
\subsection{Acceptance, adaptivity and restarts}
% Move-scaled T0 (factor x 100), time-based exponential cooling, scores
% 25/10/3/0, runtime-normalized weight updates, minimum weight, restart from
% randomized greedy + reheat after 700 rejections; Or-opt polish-and-refill
% hook (grid search disabled it in the final configuration).

% -------------------------------------------------------------------------
\section{Experimental Investigation of the Components}
\subsection{Setup and evaluation metric}
% 20 provided instances (50-1000 customers, 2-25 couriers, 3-8 skills,
% skill parameter k in {0,...,2}); pure heuristic Python, no exact-solver
% components, no third-party packages; conformance to the required command
% line (python3 skillvrp.py <instance> <timeout>) and output format; empty
% plan as always-feasible fallback; checker.py validation; %UB metric;
% seeds and noise level.
\subsection{Initial solution quality}
% Grading criterion "initial solution heuristic": multi-start greedy reaches
% 49.5% UB on average (range 43-56); no single ordering dominates ->
% best-of-four; weakness = myopic customer selection, corrected by the
% search (+20-47% per instance); also the restart fallback.
\subsection{Parameter grid search}
% 30 s runs on three representative instances x 2 seeds; searched dimensions
% (temperature factor, destroy-size profile, candidate-pool cap, ...);
% result: robust plateau, small destroy profile wins; final configuration
% table.
\subsection{Operator ablation study}
% Leave-one-out over all 23 operators, 3 instances x 2 seeds x 90 s; table
% of average/worst impact; findings: every operator contributes on average;
% ranking (worst-detour v2 and Shaw-related most valuable; skill-scarcity
% validates the problem-specific design); why the long tail is kept.
\subsection{Adaptive behavior}
% Weight-trajectory tracking per segment; final weights on n1000 (worst
% detour, regret-3, largest saving on top); convergence of best-solution
% timeline; evidence that adaptivity prices operators correctly.

% -------------------------------------------------------------------------
\section{Overall Performance}
\subsection{Comparison with the tuned LNS baseline at equal budget}
% 30 s head-to-head on the tuning instances: ALNS matches or beats the
% baseline (small table); interpretation w.r.t. framework overhead.
\subsection{Long-run results}
% 10-minute runs on all 20 instances (table): UB, greedy start, ALNS best,
% %UB, time of last improvement, iterations; 65.5% UB average; all runs
% feasible; improvement over greedy 20-47%.
\subsection{Convergence and scaling}
% Time-to-best grows with n (small instances converge < 40 s, large ones
% improve until the end); iteration counts; no systematic %UB degradation
% with n; hardness correlates with skill scarcity and courier availability;
% honest notes (single instance where the baseline stays ahead).

% -------------------------------------------------------------------------
\section{Conclusion}
% Summary of path and results; main lessons (yardstick first, throughput
% before operator richness, measure everything); future work: SISR string
% removals, time-aware insertion scoring, set-covering post-optimization
% over the collected route pool.

% -------------------------------------------------------------------------
\begin{thebibliography}{99}
% ~17 entries: Shaw 1998; Ropke & Pisinger 2006; Pisinger & Ropke 2007;
% Savelsbergh 1992; Vansteenwegen et al. 2009/2011; Labadie et al. 2012;
% Hu & Lim 2014; Gunawan et al. 2016/2017; Cappanera et al. 2011;
% Kovacs et al. 2012; Pillac et al. 2013; Cisse et al. 2017;
% Christiaens & Vanden Berghe 2020; Vidal et al. 2012/2022.
\end{thebibliography}

\end{document}
